% page 398 du fac-simile (image p-397 du scan)
% bloc 154.9 x 237.7 mm ; marge gauche 8.0 mm
\pgc{-12.700mm}{0.000mm}{
\l{\cel{3}390}
\l{}
\l{\sou{naracar}. (\sou{trans.}) Raportar pri fakto, enumerante detaloze}
\l{\cel{3}lua cirkonstanci diversa. - EFIS.}
\l{}
\l{narciso.\cel{2}(\sou{bot.)}\cel{2}Planto ek la familio \textquotedbl{}amarilidei\textquotedbl{}.}
\l{\cel{3}- L. narcissus. - DEFIRS.}
\l{}
\l{\sou{nardo}. Oleo parfumoza, quan, en la epoki antiqua, onu ex-}
\l{\cel{3}traktis de planto aromatoza. - DEFIS.}
\l{}
\l{\sou{narkotar}. (\sou{netrans}.) Esar en ta stando di stuporo o torporo}
\l{\cel{3}produktita artificale per la uzo di ula substanci (klo-}
\l{\cel{3}roformo, e c.) - DEFIS.}
\l{}
\l{\sou{narvalo}. (\sou{zool.)}\cel{2}Cetaceo di qua la maxilo havas longa den-}
\l{\cel{3}tego defensala. - L.\cel{1}\sou{monodon monoceros}. - DEFIS.}
\l{}
\l{\sou{naso}.\cel{2}Korbo ek oziero, kono-forma, a qua abutas ula reto}
\l{\cel{3}ube la fishi arivanta kaptesas. - FS.}
\l{}
\l{\sou{naskar}. (\sou{netran}s.)Ekirar la sino di la patrino. -\cel{2}IS.}
\l{}
\l{\sou{natar}.\cel{2}(\sou{netrans}.) Subtenar su ed avancar en aquo, per ula}
\l{\cel{3}movi di la membri. - EFIS.}
\l{}
\l{\sou{natro}. (\sou{kemio}).Korpo simpla, metala, qua esas elemento esen-}
\l{\cel{3}cala di sodo. -\cel{1}\sou{Simbolo kemiala}\cel{1}:\cel{2}Na. -}
\l{}
\l{\sou{naturo}. I. Ensemblo ek la kozi e forci qui existas eterne.}
\l{\cel{3}- II. La esenco, la atributi qui esas propra ye ula ento.}
\l{\cel{3}- III. La stando primitiva (di la homo), opoze a \textquotedbl{}civi-}
\l{\cel{3}lizeso\textquotedbl{}. - DEFIRS.}
\l{}
\l{\sou{naufrajar}.\cel{2}(\sou{netrans}.)Dicesas pri navo qua perisas kontre}
\l{\cel{3}rifi o brizanti, o dum lua navigo. - FIS.}
\l{}
\l{\sou{nautilo}.\cel{2}(\sou{zool.})\cel{2}Molusko cefalopoda, kun konko extera}
\l{\cel{3}parietoza. - L.\cel{1}\sou{nautilus pompilius}.\cel{2}DEFIS.}
\l{}
\l{\sou{nauzear}. (\sou{netrans.}) Tormentesar da la bezono vomar. - EFIS.}
\l{}
\l{\sou{navo}. I. Konstrukturo destinita a navigar sur od en mari ed}
\l{\cel{3}oceani. - II. Parto di kirko, de la portalo a la koreyo,}
\l{\cel{3}inter du rangi de pilastri, de koloni, qui subtenas}
\l{\cel{3}la vulto. EFIRS.}
\l{}
\l{\sou{navigar}.\cel{2}(\sou{netrans.)}\cel{3}Voyajar per irar sur od en aquo, altre}
\l{\cel{3}kam per nato. - DEFIRS.}
\l{}
\l{\sou{nayad}o.(\sou{mitol.}) Deajo infra-ranga, femina, qua prezidis la}
\l{\cel{3}fonti, riveri, e c.}
\l{}
\l{\sou{nazo.}\cel{4}Organo qua salias meze di la vizajo, rezideyo di}
\l{\cel{3}olfakto. - DEFIRS.}
}
